\documentclass[11pt]{article}

\usepackage[a4paper,margin=1in]{geometry}
\usepackage{hyperref}
\usepackage{enumitem}
\usepackage{titlesec}

\titleformat{\section}{\large\bfseries}{\thesection}{1em}{}
\titleformat{\subsection}{\normalsize\bfseries}{\thesubsection}{1em}{}

\title{KCPC-CS111 - Research References}
\author{Graphs II: Distance, Constraints, and Logical Structure}
\date{}

\begin{document}

\maketitle

\noindent
This workshop introduces shortest-path algorithms in weighted graphs, focusing on Dijkstra’s Algorithm and 0--1 BFS, and situates these techniques within their broader theoretical and historical context.

\section*{Foundational Papers}

\subsection*{[1] Dijkstra, E. W. (1959)}
\textit{A note on two problems in connexion with graphs.}  
Numerische Mathematik, 1, 269--271.  
\url{https://doi.org/10.1007/BF01386390}

\textbf{Why referenced:}  
Introduces Dijkstra’s algorithm, the primary shortest-path method taught in this workshop for graphs with non-negative edge weights.

\vspace{0.8em}

\subsection*{[2] Moore, E. F. (1959)}
\textit{The shortest path through a maze.}  
Proceedings of the International Symposium on the Theory of Switching, Harvard University Press, 285--292.

\textbf{Why referenced:}  
Establishes breadth-first search as a shortest-path algorithm in unweighted graphs, forming the conceptual foundation for both Dijkstra’s algorithm and 0--1 BFS.

\vspace{0.8em}

\subsection*{[3] Lee, C. Y. (1961)}
\textit{An algorithm for path connections and its applications.}  
IRE Transactions on Electronic Computers, EC-10(3), 346--365.  
\url{https://doi.org/10.1109/TEC.1961.5219222}

\textbf{Why referenced:}  
Demonstrates BFS-based shortest-path techniques in structured grid environments.

\section*{Structural and Data Structure Improvements}

\subsection*{[4] Dial, R. B. (1969)}
\textit{Algorithm 360: Shortest-path forest with topological ordering.}  
Communications of the ACM, 12(11), 632--633.  
\url{https://doi.org/10.1145/363269.363610}

\textbf{Why referenced:}  
Introduces bucket-based shortest-path methods for small integer weights, providing the theoretical basis for the 0--1 BFS optimisation.

\vspace{0.8em}

\subsection*{[5] Fredman, M. L., \& Tarjan, R. E. (1984)}
\textit{Fibonacci heaps and their uses in improved network optimization algorithms.}  
Journal of the ACM, 34(3), 596--615.  
\url{https://doi.org/10.1145/28869.28874}

\textbf{Why referenced:}  
Shows how priority queue design impacts the efficiency of Dijkstra’s algorithm.

\vspace{0.8em}

\subsection*{[6] Johnson, D. B. (1977)}
\textit{Efficient algorithms for shortest paths in sparse networks.}  
Journal of the ACM, 24(1), 1--13.  
\url{https://doi.org/10.1145/321992.321993}

\textbf{Why referenced:}  
Demonstrates extensions and optimisations of Dijkstra’s algorithm for sparse graphs.

\section*{Advanced and Modern Developments}

\subsection*{[7] Thorup, M. (1999)}
\textit{Undirected single-source shortest paths with positive integer weights in linear time.}  
Journal of the ACM, 46(3), 362--394.  
\url{https://doi.org/10.1145/316542.316548}

\textbf{Why referenced:}  
Proves that shortest-path problems can be solved in linear time under restricted weight conditions.

\vspace{0.8em}

\subsection*{[8] Duan, R., Mao, J., Mao, X., Shu, X., \& Yin, L. (2025--2026)}
\textit{Breaking the sorting barrier for directed single-source shortest paths.}  
Proceedings of STOC 2025.

\textit{A faster directed single-source shortest path algorithm.}  
arXiv preprint.

\textbf{Why referenced:}  
Included to show that single-source shortest path remains an active theoretical topic. 
These works revisit the fundamental complexity bounds underlying Dijkstra-style methods and demonstrate that even classical problems continue to admit structural improvements under refined models.

\section*{Summary}

Together, these papers trace the development of single-source shortest path algorithms: from the original formulations of BFS and Dijkstra’s method, through data-structure-driven optimisations, to modern complexity refinements. They situate the techniques presented in this workshop within their broader theoretical lineage.

\end{document}
